\section{Introduction}\label{art:introduction}

BFGS builds its search directions from changes in gradients, yet it often
works well when those gradients jump. Classical convex convergence theory
uses smoothness together with the evolving quasi-Newton matrices
\cite{FND_N92,FND_BN89,FND_P76}. Lewis--Overton's experiments and analyses
\cite{LO} exposed the difficulty of extending that theory: nonsmooth BFGS can
approach a minimizer while its inverse-Hessian approximations become highly
singular. What global or local convergence can ordinary full BFGS guarantee
on convex nonsmooth functions, and how can one prove it without assuming
the asymptotic behavior of its iterates or matrices?

We give three answers: deterministic counterexamples to every-start local
convergence; global and local convergence theorems for structured convex
families, including the mixed benchmark below; and a principal-minor and
comparison-energy argument that remains effective as the BFGS matrices
degenerate.

Lewis--Overton's Challenge~7.1~\cite{LO} gives this question a broader,
probabilistic setting. For locally Lipschitz semialgebraic objectives with
bounded sublevels and absolutely continuously distributed initial data, it
asks about nonbreakdown, stationary cluster points, R-linear convergence of
all function trial values, and identification of spectral subspaces under
partly smooth convergence. We study deterministic guarantees for convex
objectives.

Substantial positive theory precedes this work. Guo--Lewis extends Powell's
argument through smooth convex replacements on regions relevant to a BFGS
sequence~\cite[Theorems~3.2, 4.1 and~4.3]{GL18}. The resulting theory proves
Euclidean-norm convergence in arbitrary finite dimension
\cite[Corollary~4.2]{GL18}; for the mixed model
\[
                    F(u,v)=u^2+|v|,
\]
Corollary~4.4 gives only a subsequence approaching a point of $v=0$.
The paragraph after that corollary explicitly leaves convergence of every BFGS sequence
to $(0,0)$ open~\cite[Corollary~4.4 and following discussion]{GL18}.
Approaching the nonsmooth line does not force the smooth coordinate $u$ to
vanish. Corollary~\ref{art:mixed-benchmark} answers this whole-sequence question
for every admissible full-BFGS run, as a consequence of an arbitrary-dimensional
family theorem.

Two other advances clarify the remaining task. Xie--W\"achter's two-piece
unbounded affine model leads to an unbounded descent direction after finitely
many iterations~\preprintcite*[Theorem IV.1]{XW}.
Gebken's published 2026 introduction still identifies $u^2+|v|$ as unresolved.
His local criticality theorem for continuous smooth selections assumes an
iterate limit, repeated activity, and specified spectral and gradient-jump
behavior~\cite[Theorem~3.1 and (B1)]{BG}. Our positive results instead derive
convergence from hypotheses on the objective.

Throughout the paper, an \emph{admissible run} uses the unmodified full BFGS
update, an arbitrary $H_0\succ0$, fixed $0<c_1<c_2<1$, Armijo--weak-Wolfe
steps, and differentiable accepted points with nonzero gradients
(Definition~\ref{art:admissible}). The positive theorems concern every infinite
such run.

\paragraph{Deterministic limits.}
Prescribed-orbit constructions give smooth nonconvex failures
\cite{NEG_M04,NEG_M07}; nonconvexity also separates the examples of
Lewis--Zhang~\cite{NEG_LZ15} and Diao~\preprintcite*{NEG_DIAO26} from the
question here. Lewis--Overton's nonsmooth examples use exact search or
subgradient choices~\preprintcite*{NEG_LO08}, while Asl--Overton's failure
theorem concerns a scaled memoryless update~\cite{NEG_AO21}. These cited
examples do not supply convex local failure with the full update and the
specified search.
We construct convex objectives for which a fixed unit-first
Armijo--weak-Wolfe doubling/bisection search generates full-BFGS iterates
converging to a nonstationary point. For each local construction, one fixed
objective has failing starts arbitrarily close to its unique minimizer, all
with $H_0=I$, and every accepted or rejected query approaches that minimizer
across the family of runs. One construction is a maximum of four convex
$C^{2,1}$ branches with bounded Hessians; another is a maximum of four
uniformly strongly convex $C^{1,1}$ branches, three of which are not twice
differentiable at the minimizer (Theorems~\ref{art:local-failure}
and~\ref{art:local-failure-lower}). These refute every-start neighborhood
guarantees under either regularity regime.

\paragraph{Global and local convergence.}
We prove global value convergence for centered smooth-plus-sublinear
objectives with compact initial sublevels contained in the region where the
representation holds, and whole-iterate convergence when
the minimizer is unique (Theorem~\ref{art:centered}). This settles the mixed
benchmark above. For coercive finite convex degree-one homogeneous functions,
including coercive maxima of finitely many linear functions, every admissible
run converges to zero (Theorem~\ref{art:homogeneous}). This includes every
norm, for example $\norm{x}_1$, whose nonsmooth set for $n\ge2$ contains
coordinate hyperplanes, beyond the isolated nonsmooth point of the
Euclidean-norm case above. Local results
cover smooth convex plus finite polyhedral terms and finite polyhedral
partitions with smooth extensions (Theorem~\ref{art:polyhedral}). Moving
comparison geometry also yields local convergence for planar smooth maxima
under strict-contour and active-gradient nondegeneracy assumptions
(Corollary~\ref{art:planar-max}). The explicit three-quadratic example in
Theorem~\ref{art:nonradial-example} has curved switching sets and fails the radial
condition, so this extension has a concrete payoff beyond inward radial
modeling.

\paragraph{Analysis with degenerating matrices.}
The classical smooth analysis combines trace and determinant identities in
a matrix potential~\cite[Section~5]{FND_N92}. The uniform smooth secant bound
used there need not hold when gradients jump. We prove an inequality linking
each principal-minor degree to the next for the unmodified full-BFGS update.
This adjacent-degree control forces the arithmetic mean of the reciprocals
$1/(g_k^TH_kg_k)$ to diverge when gradients are bounded and predicted
decreases tend to zero. Comparison energies turn this restriction into convergence:
a field must retain information about the objective gap, control the gradient
changes it sees, and move with a controlled energy cost. The field, like every
comparison object in the paper, is used only in the proofs. Neither a bounded condition
number nor a prescribed spectral limit is needed
(Theorems~\ref{art:minor-hierarchy}--\ref{art:certificate}).

Sections~\ref{art:geometry}
and~\ref{art:energy} explain and establish the mechanism;
Sections~\ref{art:global} and~\ref{art:local} prove its global and local
consequences, Section~\ref{art:failure} gives the deterministic
counterexamples, and Section~\ref{art:verification} lists the remaining questions. The article contains complete proofs of these results.
An independent complete companion preserves the wider noncompact and
quantitative theory, with Lean proofs, exact correspondence and reproducible
audit evidence.
